\documentclass{article}


\title{Labwork 1}
\author{Pham Do Ngoc Minh}
\date{May 2026}

\begin{document}

\maketitle

\section{Implementation}
- 2 functions f(x) and f\_(x) to find the value of f(x) and f'(x)
\\
- For this lab, example is f(x) = x$^2$ so f'(x) = 2x
\\
- Calculate the new value of x using the gradient descent formula: x = x - R *
f’(x). Loop that 10 times
\\
- Print the value of x and f(x) each time


\section{Analyze the learning rate}
- When the learning rate r at 0.01, too small number makes the f(x) large, still around 70.
\\
- When the learning rate r at 0.1, number around this makes the diversity look clearer.
\\
- At 0.5, all 0.
\\
- At 0.99, the number too large makes the f(x) same as 0.1, but x can be negative.

\end{document}
